\section{Environment and agent specification}
\label{sec:environment-specification}

\subsection{Trainable components and task structure}

The experiments primarily train a \emph{goal selector}, rather than an end-to-end grid-navigation policy. The selector receives signals for the candidate goals together with task-irrelevant input coordinates, and produces one binary logit for choosing the left or right target. Supervised runs train this logit against the intended binary action. In the RL runs, the logit parameterizes a Bernoulli policy that is trained from the reward for the resulting outcome. The selector never receives pixels or a grid observation in the primary factorial analyses.

The project uses four environment levels, summarized below. Coordinates are zero-indexed $(\text{row},\text{column})$. At the first three levels the selector makes at most one target choice. In the fixed-fork and obstacle-grid variants, that choice is then passed as an explicit one-hot goal ID to a separately pretrained and frozen navigator. RouteWorld is different: the same selector supplies one branch bit for each reward-relevant fork, while a scripted controller executes every forced corridor move.

\begin{table}[!htbp]
\centering
\small
\begin{tabularx}{\linewidth}{@{}p{0.17\linewidth}p{0.29\linewidth}X@{}}
\toprule
Level & Environment and horizon & Learned choice and movement \\
\midrule
One-step choice & No spatial grid; horizon one. & The selector chooses $-1$ or $+1$. Reward is one iff that choice equals the intended label $Y$. This is the level used for the main factorial estimates and causal channel interventions. \\
Fixed fork & A $7\times7$ wall grid with 12 walkable cells. Start $(6,3)$; fork $(3,3)$; targets $(1,1)$ and $(1,5)$; horizon seven. & The selector chooses one target before rollout. Both shortest routes have seven primitive moves. On an expert route, the sole goal-dependent navigation decision is LEFT versus RIGHT at the fork; all other progress follows the corridor. A frozen navigator executes the route. \\
Obstacle grid & A $9\times9$ grid with requested obstacle probability $.22$ and horizon $2HW=162$. The start is sampled in rows 6--8; the left target in columns 0--3 and the right target in columns 5--8; each is at least four Manhattan steps from the start. Disconnected paths are carved open. & The selector again chooses one target. A frozen navigator handles genuine obstacle-dependent movement. Scientific summaries use the 32 maps on which that navigator was pretrained; unseen layouts are reported separately as a capability control. \\
RouteWorld & A perfect binary tree embedded in walls. Depths $D=1,2,4$ give grids $4\times5$, $7\times9$, and $13\times33$, with route lengths $4,9,27$ and $2^D$ visually identical terminal leaves. & Before rollout, the selector computes $D$ left/right choices from $D$ stage-indexed input rows. It sees padded route signals and a one-hot stage address, but no target marker, tree node, or history state. A scripted controller then executes the route. \\
\bottomrule
\end{tabularx}
\caption{The four task levels. The fixed fork is a small T-shaped corridor with one meaningful branch, whereas the obstacle environment is a conventional navigation grid with a single goal choice made upstream. RouteWorld adds repeated goal-relevant choices without learning corridor navigation.}
\label{tab:environment-levels}
\end{table}

RouteWorld does not train a policy from the physical maze observation. This observation has three flattened channels---agent position, walls, and generic terminal leaves---and does not identify the rewarded leaf or its route. The selector instead receives a fixed 44-dimensional vector: four padded proxy bits and four presence masks, sixteen padded exact-code bits and sixteen corresponding masks, and a four-way one-hot stage address. The complete route signals are repeated at each stage; only the stage address changes. The implementation evaluates all $D$ stage rows, converts their logits into a complete sequence of branch choices, and gives that sequence to the scripted corridor controller for physical verification. RouteWorld therefore measures repeated local goal selection and the accumulation of errors across forks. It does not test learned perception, online planning, a policy with memory, dependence on earlier branch outcomes, or recovery after a wrong turn.

The fixed fork's walkable layout is
\begin{center}
\begin{minipage}{0.28\linewidth}
\begin{verbatim}
#######
#L###R#
#.###.#
#..F..#
###.###
###.###
###S###
\end{verbatim}
\end{minipage}
\end{center}
where \texttt{S} is the start, \texttt{F} is the only goal-dependent fork, \texttt{L/R} are terminal targets, and \texttt{\#} denotes a wall. The two shortest action sequences are \texttt{UUULLUU} and \texttt{UUURRUU}. Reaching either target terminates the episode; physical reward is one only when it is the designated target. Collisions leave the agent in place and still consume a step.

\subsection{Selector architectures and trainable parameters}

The baseline selector input has 15 scalars: the fixed prefix $(P,P_{\mathrm{present}},R_1,\ldots,R_5)$ plus eight task-irrelevant state coordinates. Inactive $R_i$ positions are padded so that changing the parity degree does not change the input width. The default model is a non-residual MLP with two 64-unit ReLU hidden layers, biases, and one binary goal logit, for a total of 5,249 parameters. Individual experiments vary the width, depth, activation function, use of residual connections, or set of trainable parameters, as listed in the protocols below. When an experiment imposes a random-subspace limit $U$, it freezes the base network and trains exactly $U$ scalar coordinates through a fixed random count-sketch projection. The setting marked ``full'' trains every parameter.

The E5 algorithm comparison and the E10/E14 entropy studies use a matched 5,897-parameter actor. Its 17-dimensional input adds one reward-irrelevant context bit and one coordinate that records the binary path produced by the policy's preceding branch choices. It has the same two 64-unit ReLU hidden layers, followed by one goal head and eight binary heads for reward-irrelevant branch choices. All algorithm conditions retain all heads even when a head contributes no loss, keeping the actor architecture and number of trainable parameters fixed.

For actor--critic RL, sigmoid probabilities define the sampled binary actions. In E5, up to eight neutral branch actions are sampled before the final goal action; each neutral branch reconverges and cannot change success. The complete simulator horizon is 17 primitive actions---eight branch choices, eight forced merges, and the terminal goal choice---and the actor never observes a demonstration label for a neutral branch. It receives one terminal outcome, equal to one iff the final selected goal is correct. The policy loss uses the sum of sampled-action log probabilities times a normalized advantage, plus the configured entropy bonus. Evaluation is greedy.

The E5/E10/E14 critic is a separate 17-input scalar network with three width-128 residual GELU hidden transformations and 35,457 parameters. It is fully trained with squared reward-prediction loss and is excluded from every limit on trainable actor parameters. Actor and critic use AdamW. The standard setting uses 2,048 updates, a batch size of 250, a learning rate of $.003$, and zero weight decay, for a total of 512,000 presented contexts. E5 compares training every actor parameter with training exactly $U\in\{1,4,16,64,256,1024\}$ scalar coordinates. The E6 noise experiment uses the same width-64, two-hidden-layer actor design with a 16-dimensional interface and 5,833 parameters, but its fixed critic has width 256, depth three, and 136,193 parameters. E6 follows a fixed eight-epoch dataset schedule, so its update counts depend on dataset size and are specified in its protocol.

\subsection{Frozen navigation policy}

The navigator sees four flattened binary grid channels---agent position, obstacles, left target, and right target---plus a two-entry one-hot selected-goal ID. It is an MLP with two 128-unit ReLU hidden layers and four output logits for UP, RIGHT, DOWN, and LEFT. The fixed-fork instance therefore has 198 inputs and 42,500 parameters; the obstacle-grid instance has 326 inputs and 58,884 parameters. These navigator parameters are not included in any limit on trainable selector or actor parameters. The navigator is trained by supervised imitation of deterministic breadth-first-search shortest-path actions, rather than by RL, with AdamW, batch size 128, learning rate $.003$, zero weight decay, and a maximum of 250 epochs.

The fixed-fork navigator is trained on all 20 reachable goal-conditioned state--action rows and reaches 100\% action accuracy after 67 epochs. The obstacle navigator is trained on 3,908 rows from 32 generated maps and reaches 99.795\% after 29 epochs. Training stops after the requested $.99$ accuracy is sustained for three epochs; both networks are then frozen. For experiment families that enable sequential checks, each selector is evaluated with 128 fixed-fork and 128 obstacle-navigation rollouts. The obstacle-grid evaluation cycles through the 32 maps used for navigator training four times; it does not create 128 new layouts. These rollouts check that the selector's target choice can be passed to a competent controller on familiar maps. They are repeated measurements of the selector, rather than independent training runs. Held-out obstacle maps are reported as a separate capability check because the frozen navigator performs poorly on them.

\section{Shared construction, evaluation, and inference}

The original nine families use seeds
\[
\{11,23,37,41,53,67,71,83,97,101\}.
\]
Unless a family below overrides it, training, validation, and conflict sets contain 10,000, 4,000, and 10,000 rows; batch size is 250; learning rate is .003; AdamW weight decay is zero; and the global gradient norm is clipped at one. Eight task-irrelevant coordinates accompany the proxy and exact-code inputs. Requested proxy accuracies are constructed exactly rather than sampled only in expectation.

The default acquisition event requires a score of at least .9 at two consecutive logged checkpoints. Multi-goal pure control requires both behavioral agreement and directional causal control at least .9, each with a .1 margin over the other named candidates, at two checkpoints. Standalone decoder success usually uses .95 accuracy. These thresholds discretize continuous trajectories and do not establish literal phase transitions.

Causal evaluation changes one candidate channel on matched base inputs. In the parity baseline, flipping $P$ reverses the proxy recommendation, while flipping one active $R_i$ reverses the exact decoded $Y$. Reliance uses hard actions; probability and logit effects are secondary diagnostics. A probe is always interpreted as decodability, never as a goal label or proof of control by itself.

Intervals use 4,000 bootstrap resamples of training seeds at 95\% confidence. Matched conditions are differenced within seed before resampling. When specified by the analysis plan, repeated cells that differ only in task-irrelevant variables are first averaged within seed. The independent unit is always the trained initialization seed, never an evaluation row, episode, checkpoint, schedule, duplicated control, or scripted rollout. Analyses of when an event occurs preserve the checkpoint interval in which it occurs and retain runs that never reach the criterion as right-censored. Restricted means use the fixed experimental horizon. Areas under curves use linear-in-update trapezoids between directly observed checkpoints.

\section{Experiment protocols}

\subsection{E1--E5: learning order, counterexamples, capacity, and algorithms}

\textbf{E1.} A two-hidden-layer ReLU selector trains for 2,048 full-parameter update steps. The experiment crosses eight proxy accuracies $q\in\{.50,.60,.70,.80,.90,.95,.99,1\}$, five parity degrees, five widths $8,16,32,64,128$, and ten seeds. The primary statistic is the within-seed Spearman association of $q$, $k$, or capacity with $\rhoY-\rhoP$. The reported first crossings are limited to the tested grid values and should not be read as continuous sample-complexity estimates.

\textbf{E2.} At $q=.9$, separate proxy-only, exact-only, and competition models train for 4,096 steps with the same input width and parameterization. Full and head-only updates cross degrees 1--5, depths 1, 2, and 4, widths 8--128, ReLU/tanh/GELU activations, and residual connections on/off. A fixed width-128, depth-4 residual ReLU model adds random-subspace budgets $U\in\{1,4,16,64,256,1024\}$. The exact-decoding threshold is .95 final accuracy in eight of ten seeds; competition requires $\rhoY\geq.9$ in eight of ten.

\textbf{E3.} The 8,192-update training-dynamics experiment crosses five proxy accuracies, parity degrees 2--5, datasets of 1,000, 4,000, or 16,000 rows, four widths, and ten seeds, with 25 logarithmically spaced checkpoints. Acquisition and replacement require the relevant criterion to hold at two consecutive checkpoints. A reward plateau is defined as the first four-checkpoint window whose range is at most .005, whose level is close to terminal reward, and whose value exceeds a $q$-dependent floor. The prospectively specified joint prediction required the lower confidence bounds for both the simple-first and plateau-first fractions to exceed .5. The simple-first prediction passed this criterion, whereas the plateau-first prediction did not.

\textbf{E4.} Total $N=10,000$ is fixed while requested conflict counts $0,8,32,128,512,2048$, degrees 1--5, and unique-context fractions 2\%, 10\%, 50\%, and 100\% are crossed. Evaluation IDs are disjoint. The diverse-minus-concentrated contrast matches conflict count, $q$, $k$, and seed. A structured arm trains on two invertible failure transformations and tests two disjoint transformations.

\textbf{E5.} A 5,897-parameter actor with eight auxiliary heads for reward-irrelevant choices is shared by four algorithms. Clean SFT labels only the goal action. Trajectory SFT also applies a negative-log-likelihood loss to active, randomly labeled auxiliary branches. On-policy imitation samples two transitions from the current policy and queries the same minimal action oracle, while actor--critic RL uses terminal success reward. Four auxiliary-branch levels, six random-subspace limits $U\in\{1,4,16,64,256,1024\}$ plus full-parameter training, and ten seeds give 1,120 runs. Every condition receives 2,048 optimizer steps and 512,000 presented training contexts. A fixed width-128, depth-3 critic is excluded from the actor parameter limit. Full-parameter and random-subspace updates are not nested, so first crossings on this grid do not show that adding trainable coordinates causally reduces performance.

\subsection{E6--E9: perturbations, prior training, and context-dependent rules}

\textbf{E6.} This experiment crosses four algorithms; step-resampled, episode-static, state-static, and target-correlated biased perturbations; perturbations to the proxy observation, imitation label, or reward; five scales $0,.25,.5,1,2$; recurrence of $1,4,16,$ or 64 episodes per underlying state; three dataset sizes; and ten seeds:
\[
(4)(4)(3)(5)(4)(3)(10)=28,800.
\]
It fixes $q=.9$, $k=3$, a width-64 two-layer selector, batch size 256, and eight dataset epochs. Dataset sizes 2,560, 10,240, and 40,960 imply 80, 320, and 1,280 updates. Each episode repeats one contextual choice for four rows: step noise redraws per row, episode noise is shared across four rows, and state noise is fixed whenever the goal-relevant variables have the same underlying values. Biased noise adds $bY$, where $b=\min(.25,\sigma)$, while adjusting residual Gaussian variance.

The primary temporal contrast uses 7,680 positive-scale runs in which the proxy observation is perturbed, together with 1,920 matched zero-scale controls. Within every seed and combination of experimental conditions, it defines $g(N)=\rhoY(\sigma>0,N)-\rhoY(\sigma=0,N)$ and fits OLS slopes of $g(N)$ and $-|g(N)|$ against the base-two logarithm of the number of presentations before averaging the planned combinations. The scaling result must replicate separately in clean SFT and RL. Another 9,600 cells are exact negative controls because the perturbed variable does not enter the algorithm's training objective. A separate 100-run terminal-reward experiment fixes RL, reward noise, $N=10,240$, and recurrence 16 to compare step-level and episode-level perturbations when only one reward is emitted per episode. The repeated contextual rows are not sequential-policy rollouts. Increasing recurrence also reduces the diversity of underlying states, so these two factors are not separately identified. Target-correlated noise is a different intervention, rather than another degree of temporal persistence.

\textbf{E7.} Phase A trains the model to rely on the original proxy at $q_A\in\{.9,.99,1\}$ over widths 16--128, requiring reliance of at least .9 at three evaluations within 4,096 updates. Phase B resets the optimizer and runs for 1,024 updates after the original proxy is removed, decorrelated from the target at $q_B=.5$, reversed at $q_B=.4,.25,.1,0$, or replaced by a new .99-accurate proxy. Removal masks both $P$ and its presence bit. To test whether the original proxy still affects behavior, the evaluation restores it on a balanced set of conflicts. The proxy's half-life is operationally defined as the first time reliance falls below half its Phase-A value; one removal run never crosses this threshold and is retained as right-censored at the experimental horizon. The primary 840 runs are supplemented by 1,080 sensitivity runs at weight decay $10^{-5},10^{-4},10^{-3}$.

\textbf{E8.} Models in the old-history condition first receive $N_0\in\{512,2048,8192,32768\}$ examples labeled by the proxy. Compute-matched controls receive the same amount of training with balanced random labels, while no-history models skip this stage. Intended realignment either uses a fixed 8,192 examples or continues until behavior reaches a .90--.93 target band. Four final conditions run for 2,048 optimizer updates: one provides only examples on which the old and intended goals agree, one also removes the intended-goal channels, and two label disagreement cases according to the old goal 55\% or 75\% of the time. Reactivation requires old-goal reliance of at least .9 at two logged checkpoints and retains interval censoring. The complete grid contains 2,880 runs, but only 479 of 720 distinct cells meet the pre-perturbation inclusion criterion before the four final conditions are replicated. History comparisons use the compute-matched random-label control. Because inclusion depends strongly on model width, the results do not by themselves support broad scaling claims.

\textbf{E9.} Two independent rules $P_0,P_1$ and a balanced unobserved environment variable $E$ define the target $Y=P_E$. During training, the observed context satisfies $C=E$ and $C_{\mathrm{present}}=1$. Full-parameter updates cross seven widths, three depths, and ten seeds; a fixed large residual network adds six random-subspace budgets. Every selector is paired with two single-rule calibration models. The planned analysis retains only runs in which both calibration models reach .9; 38 calibration runs fail this requirement. The raw architecture heatmaps intentionally include all runs. To evaluate strict switching, the analysis takes rows on which the two rules disagree and evaluates each row twice, once with the context fixed to each possible value. Context randomization and forced mismatches preserve $E$, $P_0$, $P_1$, and the target while altering only $C$; context removal sets $C=0,C_{\mathrm{present}}=0$.

\subsection{E10--E13: entropy, alternative rules, multiple candidates, and repeated forks}

\textbf{E10.} The E5 actor--critic is trained for 2,048 updates over four proxy accuracies, three parity degrees, six entropy coefficients $0,.003,.01,.03,.1,.3$, and ten seeds, giving 720 runs. REINFORCE and a 1,024-scalar actor subspace add 120 controls for the pre-specified reference setting. Reliable control requires eight of ten seeds above .9. The decision rule specified before the experiment required an improvement of at least .20 with a positive confidence interval before running the planned mechanism panel. The largest estimate with the predicted sign was .125, so that panel was not run. E14 is a separate post-hoc response to a different setting.

\textbf{E11.} Balanced five-bit parity, majority, and conjunction, plus a six-bit multiplexer, share a six-channel padded interface. Four rules, three proxy accuracies, five widths, two depths, and ten seeds give 1,200 runs. Separate proxy, exact-rule, and competition models train for 4,096 updates. Exact-rule calibration runs repeated across $q$ are bit-identical, rather than independent decoder replications. The causal tests change the minimum possible number of input bits needed to flip each rule. Because these Boolean interventions can differ in size, their effects on output probability are not compared directly across rules.

\textbf{E12.} A direct proxy $P$, parity-coded proxy $Q$, and exact parity-coded $Y$ share one selector. The 960-run grid crosses two $q_P$ values, three $q_Q$ values, two $Q$ degrees, two $Y$ degrees, two widths, independent versus maximally nested error patterns, and ten seeds. The error rate of each proxy is fixed exactly. In the independent condition, their overlap is the nearest finite-sample value to the product of the two error rates; in the nested condition, the smaller set of errors lies wholly inside the larger set. Candidate calibration and competition models use matched masked inputs and 2,048 full-parameter update steps. Three 10,000-row evaluation sets make both proxies wrong, only $P$ wrong, or only $Q$ wrong. A model receives a pure-control label only when its closest behavioral pattern agrees with the candidate that has the largest causal influence, its error on the relevant evaluation set is at most .10, and directly flipping that candidate changes the action on at least half of cases. ``Mixture or unresolved'' is a conservative negative classification, not evidence of literal internal mixing. A post-hoc model check also attempted to reconstruct the output as a fixed additive combination of candidate rules, but this check was not part of the prospective decision rule.

\textbf{E13.} RouteWorld embeds an equal-depth binary decision tree in walls. A scripted controller executes the corridors, while the selector sees every route channel and a one-hot stage address. Depths 1, 2, and 4, three proxy accuracies, two parity degrees, two widths, two training-allocation regimes, and ten seeds give 720 runs, each with 4,000 exactly balanced route episodes. Fixed-total models receive 1,024 updates; per-fork-matched models receive $1,024D$. Evaluation includes in-distribution routes, routes on which all stages place the proxy and exact rule in conflict, and routes with one conflict stage at a time. The pre-specified approximation to whole-route success raises the pooled branch accuracy to the power $D$; a separate calculation that multiplies the measured accuracy at each stage is exploratory. The two depth-one training regimes are duplicate controls. Physical rollouts are repeated measurements, so this experiment tests repeated local selection rather than learned planning or memory.

\section{Follow-up protocols}

\textbf{E14: timing of entropy regularization.} After inspecting E10, the analysis selected the responsive $q=.75,k=4$ setting and the pre-specified $q=.9,k=3$ reference setting for six schedules: no entropy, continuous entropy coefficient $.3$, early-only entropy, delayed-only entropy, delayed entropy with an actor-optimizer reset, and delayed entropy with a critic and critic-optimizer reset. The schedule changes at update 64 of 2,048. All 120 runs use the E5 actor--critic and the ten original seeds. The primary endpoint is final intended-rule reliance on the conflict set under fixed, paired schedule comparisons. Split schedules with no effective change at the boundary reproduce uninterrupted endpoints. The analysis that groups endpoints by finite truth-table fractions, along with four deterministic replays, was selected only after the endpoints were inspected.

\textbf{E15: learning order with three candidate rules.} This follow-up was selected after inspecting the earlier results. It fixes $q_P=.9$, degree-two $Q$, width 64, depth two, and 8,192 competition updates, while crossing $q_Q\in\{.9,.95,.99\}$, $k_Y\in\{3,5\}$, independent or nested proxy errors, and 20 seeds (ten archived and ten fresh), giving 240 runs. Standalone calibration models stop at 2,048 updates. Separate 2,048-row probe-fitting and 4,096-row held-out sets balance all eight combinations of candidate-rule outputs and enumerate the raw input codewords. Probes use affine ridge fits, with permuted labels and truth tables chosen to be orthogonal to the candidate rule as controls. Equality checks on archived outputs cover the recorded shared metrics; they do not establish bit-for-bit equality of hidden states where hashes are unavailable.

\textbf{E16: adding a previously absent kind of counterexample.} This experiment fixes the nested E15 setting $q_P=.9,q_Q=.95,k_Q=2,k_Y=5$, width 64, depth two, $N=10,000$, and 8,192 updates. It inserts rows from the previously absent case in which only $Q$ is wrong, while holding the overall error rate of each proxy fixed. The seven primary counts are
\[
m\in\{0,2,10,50,100,250,450\}.
\]
Twenty paired seeds give 140 primary runs. The decision rule fixed before analyzing these counts required the prevalence of the conditional rule to drop by at least .30 overall and by at least .20 between adjacent tested values. It then selected untested quarter points from the earliest interval tied for the largest drop. The observed .55 drop between 2 and 10 inserted rows met this rule, so $m=4,6,8$ were run for 60 additional models before their outcomes were inspected. The $m=450$ condition matches the independent condition's table of error counts but is not an exact replay. Because the same inserted rows are presented repeatedly during training, $m$ counts distinct examples of this type rather than their total number of presentations.

\textbf{E17: removing the initially dominant proxy.} Twenty fresh seeds receive paired Phase-A training histories with independent or nested proxy errors for 45 updates at $q_P=q_Q=.9,k_Q=2,k_Y=3$, width 64, and depth two. Inclusion requires both resulting models to be purely controlled by $P$; all pairs pass. The difference in how accurately a probe decodes $Q$ is analyzed as a possible predictor of later behavior, but it is not used to decide which models are included. Phase B makes $P$ exactly uninformative after conditioning on the other inputs, while keeping $Q$ 90\% accurate and $Y$ exact. The 120 trajectories include independent- and nested-history models that either retain or reset optimizer state, models trained from scratch, and compute-matched controls whose earlier labels were balanced and random. The primary outcome is the area under the curve for behavioral and causal control by $Q$ through update 128. Comparisons fixed before Phase B first contrast the independent-history reset models with the nested-history reset, scratch, and random-label controls, and then compare retained versus reset optimizer state. The scratch and random-label controls do not begin with the same strong $P$ control. Consequently, only the independent-versus-nested reset comparison cleanly tests how two histories affect models that had already learned the same dominant rule.

\textbf{E19: editing first-layer weights connected to $Q$.} Five archived E17 seeds were used to select the weight columns and comparisons. A separate three-seed pilot then met the pre-specified manipulation criteria in two of three seeds, after which 20 new seeds supplied paired independent- and nested-error histories. In independent-history models, the experiment restores the two active first-layer columns connected to $Q$ exactly to their initialization values. In nested-history models, it replaces those columns with the corresponding values from the independent-history model. These edits are compared with both no edit and an equal signed displacement applied to padding columns that carry no goal information. All conditions begin Phase B under pure $P$ control with fresh optimizers. The two primary area-under-the-curve comparisons for $Q$ control require an effect of at least .02, a positive confidence interval, and a positive effect in at least 15 of 20 seeds. Control edits to padding columns must remain within .01 of the no-edit condition. The change in selective $Q$ probe accuracy must be at most $-.05$ after restoration and at least $+.05$ after transplantation, while the corresponding changes for the padding control and frozen $P/Y$ probes must remain within $[-.05,.05]$. These criteria permit the conclusion that the edited columns contribute to later control by $Q$. They do not imply that the columns form a unique circuit, explain a measured fraction of the history effect, or contain every representation of $Q$.

\textbf{E20: changing the order of otherwise identical training components.} The width-64, depth-two selector receives the literal input $(P,P_{\rm present},R_1,R_2,R_3,Q_{\rm present},Q_1,Q_2)$ over all $2^6=64$ varying codewords. The model receives no marker that identifies the training component, schedule, state, padding, row, or copy. For component $G\in\{P,Q,Y\}$, each unanimous codeword and each codeword on which $G$ is the unique minority appears twice; every other codeword appears once. This creates 96 groups, three examples from each group per minibatch, 9,216 rows, and 256 updates per component.

All six training orders begin from a bit-identical model within each seed. Model weights carry across component boundaries, while a fresh optimizer is created at every boundary. The datasets and exact ordering of individual examples are reused bit for bit whenever the same component appears; the experimental seed changes only the model initialization. If the three components are pooled, each named rule is exactly $2/3$ accurate, and every codeword on which the rules are not unanimous receives equal support for the competing labels. After 768 updates, every model receives the same 256-update agreement-only stage, containing only codewords on which all rules agree. Twenty seeds give 120 trajectories. A pilot analyzed without reference to the final outcomes learned each isolated component in all three seeds, and all 360 component blocks in the full experiment ended at the requested exact truth table.

At each checkpoint in the agreement-only stage, the primary dispersion measure averages the 15 pairwise normalized Hamming distances among the six 64-bit action tables and then integrates this value over offsets 33--128. The pre-specified criterion for persistence requires a mean area of at least .10, a lower confidence-interval endpoint above .05, and an area of at least .05 in at least 15 of 20 seeds. A signed comparison between orders that present a component last and orders that present it first tests for a recency effect. Probe and acquisition times provide secondary summaries of when representations and behavior change. The experiment measures dependence on component order under this particular construction. It does not separate label order from component-specific input frequencies, and it does not establish a universal law of phase transitions.

\section{Statistical procedures and robustness results}

\begin{table}[!htbp]
\centering
\small
\begin{adjustbox}{max width=\linewidth}
\begin{tabular}{lccccc}
\toprule
Exact-rule degree & $k=1$ & $k=2$ & $k=3$ & $k=4$ & $k=5$ \\
\midrule
Standalone decode crossing $U$ & 16 & 16 & 64 & 64 & 256 \\
Competition-control crossing $U$ & 64 & 256 & 1,024 & 1,024 & $>1,024$ \\
\bottomrule
\end{tabular}
\end{adjustbox}
\caption{In E2, these are the first tested random-subspace budgets at which the criteria are met. They are values from a discrete experimental grid, not continuous estimates of an intrinsic dimension.}
\label{tab:lw-decode-thresholds}
\end{table}

\begin{table}[!htbp]
\centering
\small
\begin{adjustbox}{max width=\linewidth}
\begin{tabular}{lccc}
\toprule
Active perturbation & Step $-$ episode & Episode $-$ state & Step $-$ state \\
\midrule
Proxy observation (four algorithms) & $+.0187\ [+.0143,+.0230]$ & $+.0991\ [+.0880,+.1116]$ & $+.1178\ [+.1065,+.1302]$ \\
Imitation label (three algorithms) & $+.0288\ [+.0254,+.0324]$ & $-.0009\ [-.0073,+.0054]$ & $+.0279\ [+.0210,+.0351]$ \\
Dense reward (RL) & $+.0062\ [+.0027,+.0100]$ & $-.0398\ [-.0484,-.0307]$ & $-.0336\ [-.0416,-.0251]$ \\
\bottomrule
\end{tabular}
\end{adjustbox}
\caption{Matched E6 contrasts in $\rhoY$ across temporal perturbation patterns. The ordering depends on where the perturbation is applied. In the 9,600 negative-control settings, the perturbed variable does not enter the training objective and behavior is exactly unchanged.}
\label{tab:lw-noise-locations}
\end{table}

\begin{table}[!htbp]
\centering
\small
\begin{tabular}{lrrr}
\toprule
E8 Stage-1 condition & Scheduled cells & Included & Excluded \\
\midrule
Fixed volume & 360 & 270 & 90 \\
Behavior matched & 360 & 209 & 151 \\
\midrule
Total & 720 & 479 & 241 \\
\bottomrule
\end{tabular}
\caption{Application of the E8 inclusion criterion before the four final conditions are replicated. The 241 excluded cells correspond to 964 run summaries after replication.}
\label{tab:lw-e8-eligibility}
\end{table}

\begin{table}[!htbp]
\centering
\small
\begin{adjustbox}{max width=\linewidth}
\begin{tabularx}{\linewidth}{l >{\raggedright\arraybackslash}X >{\raggedright\arraybackslash}X}
\toprule
Experiment & Primary result & Qualification \\
\midrule
E14 & Delayed entropy partially increased intended-rule reliance in two selected settings; endpoints usually corresponded to deterministic subsets of the finite truth table. & The settings were selected post-hoc, and the pre-specified reference-setting threshold failed; the result does not show that entropy provides a general remedy. \\
E15 & All 240 trajectories were controlled by $P$ first; common sequences included $P\to Q$, $P\to Y$, and $P\to Q\to Y$. & The grid was selected after earlier results. Probes establish decodability, but the ordering of behavioral and causal changes could not be resolved. \\
E16 & The fraction of models using the context-dependent rule fell sharply when the count rose from 2 to 10, meeting the criterion for testing counts 4, 6, and 8. & Each distinct row was repeated hundreds of times, and responses were not monotonic across seeds or at the largest count. \\
E17 & The independent-minus-nested area for early $Q$ control was $+.0596\ [+.0495,.0721]$; retaining versus resetting optimizer state changed it by $-.0018$ and $-.0100$; all models ultimately reached $Y$. & Optimizer-state effects were negligible at this scale. The two histories differ in which examples the proxies get wrong, and the probe difference is not proven to cause the later change in control. \\
E19 & Restoring the selected weights delayed $Q$ control by $+.0361\ [.0296,.0432]$ AUC; transplanting them accelerated it by $+.0293\ [.0249,.0339]$. & The result identifies a contribution from the edited input weights, not a complete circuit or a permanently necessary representation. \\
E20 & Dispersion AUC was $.39835\ [.39588,.39980]$; the last-minus-first comparison was $.99449\ [.98620,.99936]$. & Pooled training ties the competing labels, and the final finite stage contains only cases on which all rules agree, so it cannot distinguish among them. \\
\bottomrule
\end{tabularx}
\end{adjustbox}
\caption{Summary of follow-up results selected after earlier analyses. Confidence intervals resample paired training seeds; they do not correct for the earlier selection of experimental settings.}
\label{tab:lw-adaptive-ledger}
\end{table}

\FloatBarrier
\section{Extended results and robustness figures}

The remaining figures are collected here in their original visual style. Their captions state the principal result and the most important immediate qualification. The protocol descriptions and tables above define the exact comparisons.

\begin{figure}[!htbp]
  \centering
  \includegraphics[width=\linewidth]{fig11_rule_families_followup.pdf}
  \caption{When exact rules with the same number of input channels compete against a 99\%-accurate proxy, outcomes range from proxy control to intended-rule control even though standalone calibration accuracy is near one. The comparison shows that channel count does not provide a portable measure of rule complexity; it does not establish a universal ordering of Boolean or natural-language goals.}
  \label{fig:lw-rules}
\end{figure}

\begin{figure}[!htbp]
  \centering
  \includegraphics[width=\linewidth]{fig10_rl_entropy_followup.pdf}
  \caption{Across the tested entropy coefficients, the largest change in RL behavior occurs at $q=.75,k=4$, while no $q=.99$ setting becomes reliably controlled by the intended rule. The highlighted maximum is one of 24 setting-level comparisons, and the pre-specified threshold for the reference-setting mechanism analysis was not met.}
  \label{fig:lw-entropy}
\end{figure}

\begin{figure}[!htbp]
  \centering
  \includegraphics[width=\linewidth]{fig14_entropy_timing_exploratory.pdf}
  \caption{In the selected settings, introducing entropy regularization after update 64 partially increases intended-rule reliance. Endpoint values cluster near fractions of the finite parity truth table, which is consistent with nearly deterministic behavior on subsets of codewords rather than uniformly stochastic compromise. The schedule design and replay analysis are post-hoc.}
  \label{fig:lw-entropy-timing}
\end{figure}

\begin{figure}[!htbp]
  \centering
  \includegraphics[width=\linewidth]{fig5_temporal_noise.pdf}
  \caption{The effect of proxy corruption depends on its temporal structure. Corruption tied to an underlying state preserves more proxy control than freshly sampled corruption, while the behavioral difference between fresh zero-mean corruption and the unperturbed condition grows over the tested sample and update schedule.}
  \label{fig:lw-noise}
\end{figure}

\begin{figure}[!htbp]
  \centering
  \includegraphics[width=\linewidth]{fig7_hysteresis.pdf}
  \caption{After ambiguous retraining, the endpoints for old-history, compute-matched random-label, and no-history conditions nearly overlap. Longer training on the old goal nevertheless requires substantially more intended-goal data to realign behavior.}
  \label{fig:lw-hysteresis}
\end{figure}

\begin{figure}[!htbp]
  \centering
  \includegraphics[width=\linewidth]{fig8_multiplicity_capacity.pdf}
  \caption{When only an intermediate number of parameters can change, separate calibration models learn both rules successfully, but the competition model does not yet switch reliably between them according to context.}
  \label{fig:lw-multiplicity}
\end{figure}

\begin{figure}[!htbp]
  \centering
  \includegraphics[width=\linewidth]{fig12_competing_goals_followup.pdf}
  \caption{Although each of the three candidate rules can be decoded when trained in isolation, which rule controls the competition model depends on proxy accuracy, encoding, model capacity, and the overlap between the proxies' errors. Some models also use context-dependent combinations of the rules.}
  \label{fig:lw-three-goals}
\end{figure}

\clearpage
\begin{landscape}
\begin{figure}[p]
  \centering
  \includegraphics[width=0.86\linewidth,trim=0 18bp 0 25bp,clip]{fig15_multigoal_dynamics.pdf}
  \caption{Across the selected three-rule settings, all 240 training trajectories are controlled first by the direct proxy. Many later shift either to the encoded proxy or directly to the exact rule. Because these settings were selected after earlier results, the repeated learning sequences do not establish a universal law of training phases.}
  \label{fig:lw-multigoal-dynamics}
\end{figure}
\end{landscape}
\clearpage

\begin{landscape}
\begin{figure}[p]
  \centering
  \includegraphics[width=0.78\linewidth]{fig16_support_completion.pdf}
  \caption{Adding a small number of previously absent examples in which only $Q$ is wrong changes which rule controls many models. These distinct examples are sampled repeatedly during training. Transfer tests and reversals at the seed level show that the result is not a uniform switch shared by every model.}
  \label{fig:lw-support}
\end{figure}
\end{landscape}
\clearpage

\begin{landscape}
\begin{figure}[p]
  \centering
  \includegraphics[width=0.76\linewidth]{fig17_winner_knockout.pdf}
  \caption{After the direct proxy becomes uninformative, models whose earlier proxy errors were independent shift to the encoded proxy sooner than models whose errors were nested. The histories also differ in how accurately $Q$ can be decoded and in which examples each proxy gets wrong. Every condition ultimately converges to the exact rule.}
  \label{fig:lw-knockout}
\end{figure}
\end{landscape}
\clearpage

\begin{figure}[!htbp]
  \centering
  \includegraphics[width=\linewidth]{fig19_active_q_pathway.pdf}
  \caption{Restoring selected first-layer $Q$ weights to their initialization values delays early control by $Q$, while transplanting the corresponding weights from independent-error histories accelerates it. The two edits also change selective $Q$ decodability in opposite directions, whereas matched edits to padding weights have little effect. These interventions do not isolate a complete circuit.}
  \label{fig:lw-pathway}
\end{figure}

\begin{figure}[!htbp]
  \centering
  \includegraphics[width=\linewidth]{fig13_routeworld_followup.pdf}
  \caption{As RouteWorld depth increases, accuracy at each local fork decreases, and allocating more updates per fork recovers only part of this loss. The remaining local errors then reduce whole-route success. An exploratory calculation finds that multiplying stage-level accuracies closely predicts route success. Navigation between forks is scripted.}
  \label{fig:lw-routeworld}
\end{figure}

\FloatBarrier
